% © 2025 Ing. David Jaroš — CC BY-NC-ND 4.0
%
% This work is licensed under a Creative Commons Attribution-NonCommercial-NoDerivatives
% 4.0 International License (CC BY-NC-ND 4.0).
%
% License History: Earlier drafts (up to v0.3) were released under CC BY 4.0.
% From v0.4 onward, all material is released under CC BY-NC-ND 4.0 to protect
% the integrity of the theoretical work during ongoing academic development.
%
% See LICENSE.md for full license text.

\documentclass[12pt]{article}
\usepackage{amsmath,amssymb,amsthm}
\usepackage{mathtools}
\usepackage{geometry}
\usepackage{hyperref}
\geometry{a4paper, margin=1in}

% Theorem environments
\newtheorem{definition}{Definition}[section]
\newtheorem{lemma}[definition]{Lemma}
\newtheorem{theorem}[definition]{Theorem}
\newtheorem{corollary}[definition]{Corollary}
\newtheorem{proposition}[definition]{Proposition}
\newtheorem{remark}[definition]{Remark}

% Custom commands
\newcommand{\B}{\mathbb{B}}
\newcommand{\C}{\mathbb{C}}
\newcommand{\R}{\mathbb{R}}
\renewcommand{\H}{\mathbb{H}}
\newcommand{\Tr}{\mathrm{Tr}}
\newcommand{\re}{\mathrm{Re}}
\newcommand{\im}{\mathrm{Im}}
\newcommand{\calT}{\mathcal{T}}
\newcommand{\calG}{\mathcal{G}}
\newcommand{\calE}{\mathcal{E}}

\title{Appendix: General Relativity Embedding of the\\
       Unified Biquaternion Theory\\[0.3em]
       \large Hilbert Variation and the Stress-Energy Tensor $T_{\mu\nu}$}
\author{Ing.~David~Jaroš}
\date{2025}

\begin{document}

\maketitle

\noindent\textbf{License:} © 2025 Ing.~David~Jaroš. CC BY-NC-ND 4.0.
See \url{https://creativecommons.org/licenses/by-nc-nd/4.0/} for details.

\vspace{1em}

\begin{abstract}
This appendix provides a self-contained derivation of the gravitational
embedding of the Unified Biquaternion Theory (UBT) into General Relativity (GR).
Starting from the complete UBT action functional $S[\Theta, g]$, we perform
the Hilbert variation with respect to the spacetime metric $g^{\mu\nu}$ to
obtain the stress-energy tensor
\[
  T_{\mu\nu} = -\frac{2}{\sqrt{-g}}\,\frac{\delta S_{\mathrm{matter}}}{\delta g^{\mu\nu}}.
\]
We derive explicit forms for the kinetic, potential, and gauge contributions,
lift the result to the full biquaternionic tensor $\calT_{\mu\nu}$, and show
that combining with the Einstein--Hilbert term reproduces Einstein's field
equations $G_{\mu\nu} = 8\pi G\,T_{\mu\nu}$.  In the limit $\Theta \to \phi$
(real scalar field), the stress-energy tensor reduces to the classical
Klein--Gordon form, confirming GR compatibility.
\end{abstract}

\tableofcontents

\newpage

%-----------------------------------------------------------------------
\section{Action Functional $S[\Theta, g]$}
\label{sec:action}

\subsection{Conventions}

We work in Lorentzian signature $(-,+,+,+)$ throughout.
Indices are raised and lowered with $g^{\mu\nu}$ and $g_{\mu\nu}$.
We write $g := \det(g_{\mu\nu}) < 0$ and $\sqrt{-g} > 0$.
Natural units $\hbar = c = 1$ are used.

The covariant derivative acting on the biquaternionic field $\Theta$ is
\begin{equation}
  \nabla_\mu \Theta = \partial_\mu \Theta + \Omega_\mu \Theta
                      + i\,g_{\mathrm{YM}} A_\mu \Theta,
  \label{eq:cov_deriv}
\end{equation}
where $\Omega_\mu$ is the spin connection derived from the spacetime metric
$g_{\mu\nu}$ (real projection of the biquaternionic connection $\mathcal{G}_{\mu\nu}$),
and $A_\mu$ is the Yang--Mills gauge connection
$(\mathfrak{su}(3)\oplus\mathfrak{su}(2)\oplus\mathfrak{u}(1))$-valued.

\subsection{The UBT Field}

The biquaternionic field $\Theta(q,\tau) \in \B = \C\otimes\H$ is defined
over complex time $\tau = t + i\psi$ (see main text, Section~2).
Its general decomposition is
\begin{equation}
  \Theta(q,\tau) = g_{\mu\nu}(x) + i\,\psi_{\mu\nu}(x)
                  + \mathbf{j}\,\xi_{\mu\nu}(x) + \mathbf{k}\,\chi_{\mu\nu}(x),
  \label{eq:theta_decomp}
\end{equation}
where $g_{\mu\nu}$ is the real component (classical metric), and
$\psi_{\mu\nu}$, $\xi_{\mu\nu}$, $\chi_{\mu\nu}$ are imaginary/quaternionic
components encoding phase curvature.

\subsection{Complete Matter Action}

\begin{definition}[UBT matter action]
\label{def:action}
The matter sector of the UBT action, projected onto the effective
4-dimensional Lorentzian spacetime, is:
\begin{equation}
  S_{\mathrm{matter}}[\Theta, g]
  = S_{\mathrm{kin}} + S_{\mathrm{pot}} + S_{\mathrm{gauge}},
  \label{eq:S_matter}
\end{equation}
where the three terms are:
\begin{align}
  S_{\mathrm{kin}}[\Theta, g]
    &= -\frac{1}{2}\int d^4x\,\sqrt{-g}\;
       g^{\alpha\beta}\,\Tr\!\bigl[(\nabla_\alpha\Theta)^\dagger
       \nabla_\beta\Theta\bigr],
       \label{eq:S_kin} \\[4pt]
  S_{\mathrm{pot}}[\Theta, g]
    &= -\int d^4x\,\sqrt{-g}\; V(\Theta),
       \label{eq:S_pot} \\[4pt]
  S_{\mathrm{gauge}}[\Theta, g]
    &= -\frac{K}{4}\int d^4x\,\sqrt{-g}\;
       \Tr\!\bigl[F_{\alpha\beta}F^{\alpha\beta}\bigr].
       \label{eq:S_gauge}
\end{align}
Here $K > 0$ is the gauge kinetic normalisation ($K = 1$ for Abelian gauge),
$V(\Theta) = \tfrac{\lambda}{4}(\langle\Theta,\Theta\rangle - v^2)^2$
is the Mexican-hat potential, and $F_{\mu\nu} = \partial_\mu A_\nu - \partial_\nu A_\mu
+ i\,g_{\mathrm{YM}}[A_\mu, A_\nu]$ is the Yang--Mills field strength.
\end{definition}

\begin{remark}[Sign of $S_{\mathrm{kin}}$]
The minus sign in \eqref{eq:S_kin} is required by the $(-,+,+,+)$ convention:
the kinetic energy density $T^{00} = \tfrac{1}{2}(\partial_t\Theta)^2 +
\tfrac{1}{2}(\nabla\Theta)^2 > 0$ is positive with this choice.
\end{remark}

\subsection{Metric Dependence}

The metric $g^{\mu\nu}$ appears in three places in $S_{\mathrm{matter}}$:
\begin{enumerate}
  \item \textbf{Explicit:} through the index contraction $g^{\alpha\beta}$ in
        $S_{\mathrm{kin}}$ and the index raising in $F^{\alpha\beta} =
        g^{\alpha\mu}g^{\beta\nu}F_{\mu\nu}$.
  \item \textbf{Integration measure:} through the volume factor $\sqrt{-g}$
        in every term.
  \item \textbf{Covariant derivative:} through the spin connection $\Omega_\mu$,
        which is determined by $g_{\mu\nu}$ via the tetrad.  On-shell (when
        $\Theta$ satisfies its Euler--Lagrange equations), the spin-connection
        variation contributes only boundary terms that vanish for variations
        with compact support.
\end{enumerate}

%-----------------------------------------------------------------------
\section{Hilbert Variation of the Matter Action}
\label{sec:variation}

\begin{definition}[Hilbert stress-energy tensor]
\label{def:hilbert}
Given the matter action $S_{\mathrm{matter}}[\Theta, g]$, the
\emph{Hilbert stress-energy tensor} is defined by:
\begin{equation}
  \label{eq:hilbert}
  T_{\mu\nu} := -\frac{2}{\sqrt{-g}}\,
    \frac{\delta S_{\mathrm{matter}}}{\delta g^{\mu\nu}}.
\end{equation}
The sign convention ensures that $G_{\mu\nu} = 8\pi G\,T_{\mu\nu}$ holds on
the right-hand side of Einstein's equations.
\end{definition}

\subsection{Standard Variational Identities}

\begin{lemma}[Metric variation identities]
\label{lem:var_ids}
Under $g^{\mu\nu} \to g^{\mu\nu} + \delta g^{\mu\nu}$:
\begin{align}
  \delta\sqrt{-g}
    &= -\tfrac{1}{2}\sqrt{-g}\; g_{\mu\nu}\,\delta g^{\mu\nu},
    \label{eq:var_sqrtg} \\
  \delta g_{\alpha\beta}
    &= -g_{\alpha\mu}\,g_{\beta\nu}\,\delta g^{\mu\nu},
    \label{eq:var_lower} \\
  \delta\!\bigl(\sqrt{-g}\,g^{\alpha\beta}\bigr)
    &= \sqrt{-g}\Bigl[\delta^{(\alpha}_\mu\delta^{\beta)}_\nu
       - \tfrac{1}{2}g_{\mu\nu}\,g^{\alpha\beta}\Bigr]\delta g^{\mu\nu}.
    \label{eq:var_combined}
\end{align}
\end{lemma}

\begin{proof}
Identity~\eqref{eq:var_sqrtg}: from $\delta\ln|\det g| = g^{\mu\nu}\delta g_{\mu\nu}
= -g_{\mu\nu}\delta g^{\mu\nu}$ (using $g^{\mu\alpha}g_{\alpha\nu}=\delta^\mu_\nu$).
Identity~\eqref{eq:var_lower}: differentiate $g^{\mu\alpha}g_{\alpha\nu}=\delta^\mu_\nu$.
Identity~\eqref{eq:var_combined}: combine the first two with the product rule
and symmetrise in $(\mu,\nu)$.
\end{proof}

\subsection{Kinetic Term Variation}

Write $S_{\mathrm{kin}} = -\tfrac{1}{2}\int d^4x\,\sqrt{-g}\,
g^{\alpha\beta}\,K_{\alpha\beta}$, where
$K_{\alpha\beta} := \Tr[(\nabla_\alpha\Theta)^\dagger\nabla_\beta\Theta]$.

\begin{proposition}[Kinetic contribution to $T_{\mu\nu}$]
\label{prop:T_kin}
For $\Theta$ satisfying the UBT Euler--Lagrange equations, on-shell:
\begin{equation}
  \label{eq:T_kin}
  \boxed{
    T^{(\mathrm{kin})}_{\mu\nu}
    = \Tr\!\bigl[(\nabla_\mu\Theta)^\dagger\nabla_\nu\Theta\bigr]
      - \tfrac{1}{2}\,g_{\mu\nu}\,
        \Tr\!\bigl[(\nabla_\alpha\Theta)^\dagger\nabla^\alpha\Theta\bigr].
  }
\end{equation}
\end{proposition}

\begin{proof}
Vary $S_{\mathrm{kin}}$ with respect to $g^{\mu\nu}$:
\begin{equation}
  \delta S_{\mathrm{kin}}
  = -\frac{1}{2}\int d^4x\,\Bigl[
      \delta(\sqrt{-g}\,g^{\alpha\beta})\,K_{\alpha\beta}
    + \sqrt{-g}\,g^{\alpha\beta}\,\delta K_{\alpha\beta}
    \Bigr].
\end{equation}

\textbf{First piece (explicit metric dependence).}
Applying identity~\eqref{eq:var_combined}:
\begin{equation}
  \delta(\sqrt{-g}\,g^{\alpha\beta})\,K_{\alpha\beta}
  = \sqrt{-g}\Bigl[K_{\mu\nu} - \tfrac{1}{2}g_{\mu\nu}\,
    g^{\alpha\beta}K_{\alpha\beta}\Bigr]\delta g^{\mu\nu}.
\end{equation}

\textbf{Second piece (connection variation).}
The spin connection $\Omega_\mu$ depends on $g_{\mu\nu}$ through the tetrad.
Its metric variation produces terms $\propto (\nabla^\mu\nabla_\mu\Theta)\,\delta\Theta$
upon integration by parts.  When $\Theta$ satisfies its Euler--Lagrange equation
$\nabla^\mu\nabla_\mu\Theta = \partial V/\partial\Theta^\dagger$, these reduce to
boundary terms that vanish for compactly supported variations.
We discard this piece on-shell.

Applying Definition~\ref{def:hilbert} to the first piece:
\begin{equation}
  T^{(\mathrm{kin})}_{\mu\nu}
  = -\frac{2}{\sqrt{-g}}\cdot\bigl(-\tfrac{1}{2}\bigr)\cdot
    \tfrac{1}{2}\sqrt{-g}\bigl[K_{\mu\nu} - \tfrac{1}{2}g_{\mu\nu}K^\alpha{}_\alpha\bigr]
  = K_{\mu\nu} - \tfrac{1}{2}g_{\mu\nu}K^\alpha{}_\alpha,
\end{equation}
which is \eqref{eq:T_kin}.
\end{proof}

\subsection{Potential Term Variation}

\begin{proposition}[Potential contribution to $T_{\mu\nu}$]
\label{prop:T_pot}
The potential $V(\Theta)$ is a scalar built from $\Theta$ with no explicit
metric contractions; it depends on $g^{\mu\nu}$ only through $\sqrt{-g}$.
Therefore:
\begin{equation}
  \label{eq:T_pot}
  \boxed{T^{(\mathrm{pot})}_{\mu\nu} = -g_{\mu\nu}\,V(\Theta).}
\end{equation}
\end{proposition}

\begin{proof}
\begin{equation}
  \frac{\delta S_{\mathrm{pot}}}{\delta g^{\mu\nu}}
  = -\int d^4x\,\frac{\delta\sqrt{-g}}{\delta g^{\mu\nu}}\,V(\Theta)
  = \frac{1}{2}\sqrt{-g}\;g_{\mu\nu}\,V(\Theta).
\end{equation}
Applying Definition~\ref{def:hilbert}: $T^{(\mathrm{pot})}_{\mu\nu} = -g_{\mu\nu}V(\Theta)$.
\end{proof}

\begin{remark}[Cosmological constant form]
$T^{(\mathrm{pot})}_{\mu\nu} = -g_{\mu\nu}V$ has the structure of a
(field-dependent) cosmological constant.  At the vacuum $|\Theta|=v$,
the Mexican-hat potential vanishes: $V|_{|\Theta|=v}=0$, so
$T^{(\mathrm{pot})}_{\mu\nu}|_{\mathrm{vac}}=0$.
\end{remark}

\subsection{Gauge Kinetic Term Variation}

With $F^{\mu\nu} = g^{\mu\alpha}g^{\nu\beta}F_{\alpha\beta}$, the metric
appears in $\sqrt{-g}$ and in the two index-raising factors.

\begin{proposition}[Gauge kinetic contribution to $T_{\mu\nu}$]
\label{prop:T_gauge}
\begin{equation}
  \label{eq:T_gauge}
  \boxed{
    T^{(\mathrm{gauge})}_{\mu\nu}
    = K\Bigl[
        \Tr\!\bigl[F_{\mu\alpha}F_\nu{}^{\alpha}\bigr]
        - \tfrac{1}{4}\,g_{\mu\nu}\,\Tr\!\bigl[F_{\alpha\beta}F^{\alpha\beta}\bigr]
      \Bigr].
  }
\end{equation}
\end{proposition}

\begin{proof}
Vary $S_{\mathrm{gauge}} = -\tfrac{K}{4}\int d^4x\,\sqrt{-g}\,
\Tr[F_{\mu\nu}F^{\mu\nu}]$ term by term with respect to $g^{\rho\sigma}$:

\textit{Variation of $\sqrt{-g}$:}
$\displaystyle-\frac{K}{4}\bigl(-\tfrac{1}{2}\sqrt{-g}\,g_{\rho\sigma}\bigr)
F_{\mu\nu}F^{\mu\nu}
= \frac{K}{8}\sqrt{-g}\;g_{\rho\sigma}\Tr[F_{\mu\nu}F^{\mu\nu}]$.

\textit{Variation of the first $g^{\mu\alpha}$:}
By the antisymmetry $F_{\mu\nu} = -F_{\nu\mu}$ and the identity
$\delta g^{\mu\alpha} = -\delta^{(\mu}_\rho\delta^{\alpha)}_\sigma\,\delta g^{\rho\sigma}$:
$\displaystyle-\frac{K}{4}\sqrt{-g}\;\Tr[F_{\rho\beta}F_\sigma{}^\beta
+ F_{\sigma\beta}F_\rho{}^\beta]$.

\textit{Variation of the second $g^{\nu\beta}$:}
Identical to the previous piece by antisymmetry.

Summing all pieces and applying Definition~\ref{def:hilbert}:
\begin{equation}
  \frac{\delta S_{\mathrm{gauge}}}{\delta g^{\rho\sigma}}
  = \frac{\sqrt{-g}\,K}{4}\Bigl[
      \tfrac{1}{2}g_{\rho\sigma}\Tr[F_{\mu\nu}F^{\mu\nu}]
      - 2\Tr[F_{\rho\alpha}F_\sigma{}^{\alpha}]
    \Bigr],
\end{equation}
which gives $T^{(\mathrm{gauge})}_{\rho\sigma}$ as in \eqref{eq:T_gauge}.
\end{proof}

%-----------------------------------------------------------------------
\section{Complete Stress-Energy Tensor $T_{\mu\nu}$}
\label{sec:T_total}

\begin{theorem}[UBT matter stress-energy tensor]
\label{thm:T_matter}
The complete matter stress-energy tensor obtained by the Hilbert prescription
from $S_{\mathrm{matter}} = S_{\mathrm{kin}} + S_{\mathrm{pot}} + S_{\mathrm{gauge}}$ is:
\begin{equation}
  \label{eq:T_total}
  \boxed{
    T_{\mu\nu}
    = \underbrace{\Tr[(\nabla_\mu\Theta)^\dagger\nabla_\nu\Theta]
        - \tfrac{1}{2}g_{\mu\nu}\Tr[(\nabla_\alpha\Theta)^\dagger\nabla^\alpha\Theta]}_{
        T^{(\mathrm{kin})}_{\mu\nu}}
    + \underbrace{-g_{\mu\nu}V(\Theta)}_{
        T^{(\mathrm{pot})}_{\mu\nu}}
    + K\underbrace{\Bigl[\Tr[F_{\mu\alpha}F_\nu{}^\alpha]
        - \tfrac{1}{4}g_{\mu\nu}\Tr[F_{\alpha\beta}F^{\alpha\beta}]\Bigr]}_{
        T^{(\mathrm{gauge})}_{\mu\nu}/K}.
  }
\end{equation}
This is valid on-shell, i.e.\ when $\Theta$ satisfies the UBT Euler--Lagrange
equations $\nabla^\dagger\nabla\Theta = \partial V/\partial\Theta^\dagger$.
\end{theorem}

\begin{proof}
Additivity of the variational derivative:
$\delta S_{\mathrm{matter}}/\delta g^{\mu\nu}
= \delta S_{\mathrm{kin}}/\delta g^{\mu\nu}
+ \delta S_{\mathrm{pot}}/\delta g^{\mu\nu}
+ \delta S_{\mathrm{gauge}}/\delta g^{\mu\nu}$.
Applying Definition~\ref{def:hilbert} to each piece and summing the results
of Propositions~\ref{prop:T_kin}, \ref{prop:T_pot}, and~\ref{prop:T_gauge}
gives \eqref{eq:T_total}.
\end{proof}

%-----------------------------------------------------------------------
\section{Biquaternionic Lift to $\mathcal{T}_{\mu\nu}$}
\label{sec:calT}

Because the full UBT action $S[\Theta]$ is biquaternion-valued, the Hilbert
prescription~\eqref{eq:hilbert} yields a biquaternion-valued tensor
$\calT_{\mu\nu} \in \B$.

\begin{definition}[Biquaternionic stress-energy tensor]
\label{def:calT}
\begin{equation}
  \label{eq:calT_def}
  \calT_{\mu\nu}
  := -\frac{2}{\sqrt{-g}}\,\frac{\delta S[\Theta]}{\delta g^{\mu\nu}}
  \;\in\; \B,
\end{equation}
where $S[\Theta]$ is the complete biquaternion-valued matter action.
\end{definition}

The decomposition mirrors Theorem~\ref{thm:T_matter}:
\begin{align}
  \calT^{(\mathrm{kin})}_{\mu\nu}
    &= \Tr_{\B}\!\bigl[(\nabla_\mu\Theta)^\dagger\nabla_\nu\Theta\bigr]
       - \tfrac{1}{2}g_{\mu\nu}
         \Tr_{\B}\!\bigl[(\nabla_\alpha\Theta)^\dagger\nabla^\alpha\Theta\bigr],
       \label{eq:calT_kin} \\
  \calT^{(\mathrm{pot})}_{\mu\nu}
    &= -g_{\mu\nu}\,V(\Theta),
       \label{eq:calT_pot} \\
  \calT^{(\mathrm{gauge})}_{\mu\nu}
    &= K\Bigl[
         \Tr_{\B}\!\bigl[F_{\mu\alpha}F_\nu{}^\alpha\bigr]
         - \tfrac{1}{4}g_{\mu\nu}\Tr_{\B}\!\bigl[F_{\alpha\beta}F^{\alpha\beta}\bigr]
       \Bigr],
       \label{eq:calT_gauge}
\end{align}
where $\Tr_{\B}$ denotes the biquaternionic trace over all internal indices.

\begin{definition}[Physical stress-energy tensor]
The \emph{physical (real) stress-energy tensor} is the Hermitian projection:
\begin{equation}
  \label{eq:T_real_proj}
  T_{\mu\nu} := \re(\calT_{\mu\nu}),
\end{equation}
applied component-wise.  The imaginary part $\im(\calT_{\mu\nu})$ encodes
phase curvature and does not couple to classical matter.
\end{definition}

The fundamental UBT field equation is:
\begin{equation}
  \calE_{\mu\nu} = \kappa\,\calT_{\mu\nu},
  \quad
  \calE_{\mu\nu},\,\calT_{\mu\nu} \in \B,
\end{equation}
and the Einstein equations arise \emph{only} after taking the real projection:
\begin{equation}
  G_{\mu\nu} := \re(\calE_{\mu\nu}) = 8\pi G\,T_{\mu\nu} := 8\pi G\,\re(\calT_{\mu\nu}).
\end{equation}

%-----------------------------------------------------------------------
\section{Einstein Field Equations with Matter}
\label{sec:einstein}

\begin{definition}[Einstein--Hilbert action]
\label{def:EH}
\begin{equation}
  \label{eq:S_EH}
  S_{\mathrm{EH}}[g] = \frac{1}{16\pi G}\int d^4x\,\sqrt{-g}\; R,
\end{equation}
where $R = g^{\mu\nu}R_{\mu\nu}$ is the Ricci scalar.
\end{definition}

\begin{theorem}[Einstein field equations from UBT]
\label{thm:einstein}
The stationary condition $\delta S_{\mathrm{total}}/\delta g^{\mu\nu} = 0$
for the combined action
\begin{equation}
  S_{\mathrm{total}}[g,\Theta]
  = S_{\mathrm{EH}}[g] + S_{\mathrm{matter}}[\Theta, g]
\end{equation}
yields:
\begin{equation}
  \label{eq:einstein_matter}
  \boxed{G_{\mu\nu} = 8\pi G\,T_{\mu\nu},}
\end{equation}
where $G_{\mu\nu} = R_{\mu\nu} - \tfrac{1}{2}g_{\mu\nu}R$ is the Einstein
tensor and $T_{\mu\nu}$ is given by \eqref{eq:T_total}.
\end{theorem}

\begin{proof}
The variation of $S_{\mathrm{EH}}$ gives the standard Palatini result:
\begin{equation}
  \frac{\delta S_{\mathrm{EH}}}{\delta g^{\mu\nu}}
  = \frac{\sqrt{-g}}{16\pi G}\,G_{\mu\nu}.
\end{equation}
The variation of $S_{\mathrm{matter}}$ gives (by Definition~\ref{def:hilbert}):
\begin{equation}
  \frac{\delta S_{\mathrm{matter}}}{\delta g^{\mu\nu}}
  = -\frac{\sqrt{-g}}{2}\,T_{\mu\nu}.
\end{equation}
Setting the total variation to zero:
\begin{equation}
  \frac{\sqrt{-g}}{16\pi G}\,G_{\mu\nu}
  - \frac{\sqrt{-g}}{2}\,T_{\mu\nu} = 0
  \quad\Longrightarrow\quad
  G_{\mu\nu} = 8\pi G\,T_{\mu\nu}.
\end{equation}
\end{proof}

\begin{corollary}[Energy-momentum conservation]
\label{cor:conservation}
The contracted Bianchi identity $\nabla^\mu G_{\mu\nu} \equiv 0$ and the
Einstein equations~\eqref{eq:einstein_matter} imply automatically:
\begin{equation}
  \nabla^\mu T_{\mu\nu} = 0.
\end{equation}
\end{corollary}

\begin{proof}
$8\pi G\,\nabla^\mu T_{\mu\nu} = \nabla^\mu G_{\mu\nu} = 0$ (contracted Bianchi).
Since $8\pi G \neq 0$, the result follows.
\end{proof}

%-----------------------------------------------------------------------
\section{Covariance and Consistency Checks}
\label{sec:checks}

\subsection{Coordinate Covariance}

$T_{\mu\nu}$ defined by~\eqref{eq:hilbert} is manifestly covariant: the
variation $\delta S/\delta g^{\mu\nu}$ transforms as a $(0,2)$ tensor under
diffeomorphisms by construction.  The factor $1/\sqrt{-g}$ converts the
tensor density to a proper tensor.

\subsection{Dimensional Consistency}

In natural units ($\hbar = c = 1$):
\begin{itemize}
  \item $[\Theta] = M^{3/2}$ (mass dimension $3/2$, like a Dirac field),
  \item $[\nabla_\mu\Theta] = M^{5/2}$,
  \item $[(\nabla_\mu\Theta)^\dagger\nabla_\nu\Theta] = M^5$,
  \item $[\sqrt{-g}\,d^4x] = M^{-4}$,
  \item $[S_{\mathrm{kin}}] = M^0$ (dimensionless action),
  \item $[T_{\mu\nu}] = [1/\sqrt{-g}]\cdot[\delta S/\delta g^{\mu\nu}]
                       = M^4\cdot M^{-4} = M^4$
                       (energy density: $\mathrm{Energy/Volume}$). $\checkmark$
\end{itemize}

\subsection{Symmetry of $T_{\mu\nu}$}

Each piece is manifestly symmetric:
\begin{itemize}
  \item $T^{(\mathrm{kin})}_{\mu\nu}$: symmetric because $K_{\mu\nu} = K_{\nu\mu}$
        ($\Tr[A^\dagger B] = \overline{\Tr[B^\dagger A]}$ and the real part is symmetric).
  \item $T^{(\mathrm{pot})}_{\mu\nu} = -g_{\mu\nu}V$: symmetric by $g_{\mu\nu} = g_{\nu\mu}$.
  \item $T^{(\mathrm{gauge})}_{\mu\nu}$: symmetric by the explicit symmetrisation
        in $(\mu,\nu)$ and $F_{\mu\alpha}F_\nu{}^\alpha + F_{\nu\alpha}F_\mu{}^\alpha = 2F_{\mu\alpha}F_\nu{}^\alpha$ (due to taking the symmetrised variation $\delta g^{(\mu\nu)}$).
\end{itemize}

%-----------------------------------------------------------------------
\section{Classical Scalar-Field Limit}
\label{sec:scalar_limit}

To verify compatibility with standard field theory, we take
$\Theta \to \phi$ (real scalar field with $[\phi] = M$), no gauge fields
($A_\mu = 0$), and $\Tr[\cdot] \to \tfrac{1}{2}(\cdot)$ (normalisation).

\subsection{Scalar Field Action}

The UBT matter action reduces to:
\begin{equation}
  S_{\mathrm{scalar}} = -\frac{1}{2}\int d^4x\,\sqrt{-g}\;
    \bigl[g^{\alpha\beta}\partial_\alpha\phi\,\partial_\beta\phi + m^2\phi^2\bigr],
\end{equation}
i.e.\ the Klein--Gordon action (the potential $V = \tfrac{1}{2}m^2\phi^2$
at leading order near $\phi=0$).

\subsection{Classical Stress-Energy Tensor}

Applying Theorem~\ref{thm:T_matter} in this limit:

\begin{corollary}[Klein--Gordon stress-energy tensor]
\label{cor:kg}
For a real scalar field $\phi$ with mass $m$:
\begin{equation}
  \label{eq:T_KG}
  \boxed{
    T_{\mu\nu}^{(\mathrm{KG})}
    = \partial_\mu\phi\,\partial_\nu\phi
      - \frac{1}{2}\,g_{\mu\nu}\bigl[(\partial\phi)^2 + m^2\phi^2\bigr].
  }
\end{equation}
\end{corollary}

\begin{proof}
From \eqref{eq:T_kin} with $\Tr[(\nabla_\mu\phi)^\dagger\nabla_\nu\phi]
= \partial_\mu\phi\,\partial_\nu\phi$ (real scalar),
and from \eqref{eq:T_pot} with $V(\phi) = \tfrac{1}{2}m^2\phi^2$:
\[
  T_{\mu\nu}
  = \partial_\mu\phi\,\partial_\nu\phi
    - \tfrac{1}{2}g_{\mu\nu}(\partial\phi)^2
    - \tfrac{1}{2}m^2 g_{\mu\nu}\phi^2,
\]
which is \eqref{eq:T_KG}. $\checkmark$
\end{proof}

\begin{remark}[Massless limit]
For $m = 0$: $T^{(\mathrm{KG})}_{\mu\nu} = \partial_\mu\phi\,\partial_\nu\phi
- \tfrac{1}{2}g_{\mu\nu}(\partial\phi)^2$.
Taking the trace: $g^{\mu\nu}T_{\mu\nu} = (\partial\phi)^2 - 2(\partial\phi)^2 = -(\partial\phi)^2 \neq 0$
in general (massless scalar is not conformally invariant in $d=4$,
unlike electromagnetism).
\end{remark}

%-----------------------------------------------------------------------
\section{Reduction to General Relativity}
\label{sec:gr_reduction}

\begin{theorem}[GR embedding of UBT]
\label{thm:gr_embedding}
In the real-valued limit $\im(\calT_{\mu\nu}) \to 0$ (equivalently,
the imaginary components of $\Theta$ vanish: $\psi_{\mu\nu} = \xi_{\mu\nu}
= \chi_{\mu\nu} = 0$ in decomposition~\eqref{eq:theta_decomp}), the
biquaternionic field equations reduce to:
\begin{equation}
  G_{\mu\nu} = 8\pi G\, T_{\mu\nu},
\end{equation}
i.e.\ Einstein's field equations with the classical (real-valued) stress-energy
tensor \eqref{eq:T_total}.  This holds for arbitrary spacetime curvature,
including:
\begin{itemize}
  \item Flat spacetime ($R_{\mu\nu} = 0$): Minkowski, $T_{\mu\nu} = 0$.
  \item Weak-field limit: linearised gravity.
  \item Strong-field regimes: black holes, neutron stars, gravitational waves.
  \item Cosmological solutions: FLRW metrics with $R \neq 0$.
\end{itemize}
\end{theorem}

\begin{proof}
Theorem~\ref{thm:einstein} gives $G_{\mu\nu} = 8\pi G\,T_{\mu\nu}$ for all
metric field configurations.  In the real-valued limit, $T_{\mu\nu} = \re(\calT_{\mu\nu})
= \calT_{\mu\nu}$ (since imaginary part vanishes), so the classical Einstein
equations are recovered exactly with no correction terms.
\end{proof}

\begin{remark}[Covariance]
All curvature quantities, covariant derivatives, and stress-energy components
are defined at the biquaternionic level; GR arises as the real sector.
Coordinate transformations $x^\mu \to x'^\mu$ act identically on both sectors,
so covariance is manifest.
\end{remark}

%-----------------------------------------------------------------------
\section{Summary}
\label{sec:summary}

\begin{enumerate}
  \item \textbf{Action:} $S_{\mathrm{matter}}[\Theta,g]
        = S_{\mathrm{kin}} + S_{\mathrm{pot}} + S_{\mathrm{gauge}}$,
        Eq.~\eqref{eq:S_matter}.
  \item \textbf{Metric dependence:} $g^{\mu\nu}$ appears explicitly, in
        $\sqrt{-g}$, and through the spin connection.
  \item \textbf{Hilbert variation:}
        $T_{\mu\nu} = -(2/\sqrt{-g})\,\delta S_{\mathrm{matter}}/\delta g^{\mu\nu}$,
        Definition~\ref{def:hilbert}.
  \item \textbf{Stress-energy tensor:}
        $T_{\mu\nu} = T^{(\mathrm{kin})}_{\mu\nu}
                    + T^{(\mathrm{pot})}_{\mu\nu}
                    + T^{(\mathrm{gauge})}_{\mu\nu}$,
        Theorem~\ref{thm:T_matter}, Eq.~\eqref{eq:T_total}.
  \item \textbf{Einstein equations:} $G_{\mu\nu} = 8\pi G\,T_{\mu\nu}$
        from $\delta S_{\mathrm{total}}/\delta g^{\mu\nu} = 0$,
        Theorem~\ref{thm:einstein}.
  \item \textbf{Classical limit:} reduces to Klein--Gordon
        $T_{\mu\nu}^{(\mathrm{KG})}$
        for $\Theta\to\phi$ (real scalar), Corollary~\ref{cor:kg}.
  \item \textbf{Covariance:} manifest by construction.
  \item \textbf{Dimensional consistency:} $[T_{\mu\nu}] = M^4$. $\checkmark$
  \item \textbf{Symmetry:} $T_{\mu\nu} = T_{\nu\mu}$. $\checkmark$
  \item \textbf{Conservation:} $\nabla^\mu T_{\mu\nu} = 0$ (Bianchi identity),
        Corollary~\ref{cor:conservation}.
\end{enumerate}

\begin{thebibliography}{99}
\bibitem{wald}
  Wald, R.~M. (1984). \textit{General Relativity}. University of Chicago Press.
\bibitem{carroll}
  Carroll, S.~M. (2004). \textit{Spacetime and Geometry}. Addison--Wesley.
\bibitem{misner}
  Misner, C.~W., Thorne, K.~S., \& Wheeler, J.~A. (1973). \textit{Gravitation}. Freeman.
\bibitem{ps}
  Peskin, M.~E.\ \& Schroeder, D.~V. (1995).
  \textit{An Introduction to Quantum Field Theory}. Westview Press.
\bibitem{appendix_AA}
  Jaroš, D. (2025).
  UBT Action Principle: \texttt{consolidation\_project/appendix\_AA\_theta\_action.tex}.
\bibitem{appendix_R}
  Jaroš, D. (2025).
  GR Equivalence: \texttt{consolidation\_project/appendix\_R\_GR\_equivalence.tex}.
\end{thebibliography}

\vspace{2em}
\noindent\textbf{License:} © 2025 Ing.~David~Jaroš — CC BY-NC-ND 4.0.

\end{document}
